% !TeX program = xelatex
% !TeX encoding = utf8
% !TeX root = ProjectPlan.tex
% vim: set ts=2 sw=2 noet:

\documentclass[a4paper, twosided, 11pt]{scrartcl}

%% Language configuration
\usepackage{polyglossia}
\setdefaultlanguage{english}

%$ Specify where floatings go and make nicer looking tables
\usepackage{float}
\usepackage{array}
\usepackage{booktabs}
\usepackage{tabularx}

%% graphics
\usepackage{graphicx}

%% Set up font
\usepackage[T1]{fontenc}
\usepackage[usefilenames, DefaultFeatures={Ligatures=Common}]{plex-otf}
\renewcommand*\familydefault{\sfdefault}

%% Header and Footers
\usepackage[automark]{scrlayer-scrpage}
\ihead{Project Plan}
\chead{}
\ohead{\headmark}

\ifoot{}
\cfoot{}
\ofoot{\pagemark}

%% Access metadata
\usepackage{titling}

%% gantt configuration
\usepackage{pgfgantt}
\usepackage[]{geometry}


%% Landscape pages
\usepackage{pdflscape}

% Metadata
\title{Project Plan}
\author{Naoki Pross, Sara Cinzia Halter}


% Document
\begin{document}
\begin{titlepage}
	\includegraphics[height=3cm]{fig/ost-logo}
	\begin{flushright}
		% TODO: OST logo
		\vspace{5cm}
		{\Huge \bfseries \thetitle} \\
		\vspace{5mm}
		{\LARGE Project: \textit{Fading illustration with SDR}} \\
		\vspace{5mm}
		{\LARGE \bfseries Version 0.1}
	\end{flushright}
\end{titlepage}

\section*{Changelog}
\begin{tabularx}{\textwidth}{lrlX}
	\toprule
	\bfseries Date & \bfseries Version & \bfseries Author & \bfseries Description \\
	\midrule
	2021-09-21 & 0.1 & Pross N.  & Created document \\
	\bottomrule
\end{tabularx}

\vfill
{
	\renewcommand{\arraystretch}{2}
	\begin{tabularx}{\textwidth}{lp{.2\textwidth}X}
		                   & \bfseries Date & \bfseries Signature\\
		Halter Sara Cinzia & \hrulefill & \hrulefill \\
		Pross Naoki        & \hrulefill & \hrulefill \\
	\end{tabularx}
}

\clearpage
\tableofcontents
\listoffigures
\listoftables
\clearpage

\section{Introduction}

\section{Task Description}

\begin{itemize}
	\item Prototype of QAM transmission with SDR (add other modulations later)
	\item Model of a fading channel and simulation
	\begin{itemize}
		\item Slow fading (shadowing)
		\item Fast fading (flat fading and multi path 
	\end{itemize}
	\item Demonstration (hardware) with fading channel
	\item Demonstration with moving receiver 
\end{itemize}



\clearpage
\newgeometry{vmargin = 2.5cm, hmargin = 2cm} 

\begin{landscape}
	\section{Project Schedule}
	\begin{figure}[h]
		\begin{ganttchart}[
			time slot format = isodate,
			% style changes
			% hgrid,
			vgrid,
			x unit = 2.5mm,
			y unit chart = 1.5cm,
			bar height = .4,
			inline,
			milestone inline label node/.append style = {
				above = 4mm, anchor = south east,
				fill = white,
			},
		]{2021-09-20}{2021-12-26}

			\gantttitlecalendar{year, month = name, week = 38} \\

			\ganttgroup{Research}{2021-09-21}{2021-11-28} \ganttnewline
			\ganttbar{Theory}{2021-09-22}{2021-11-28} \ganttnewline
			\ganttbar{Software Toolboxes}{2021-09-27}{2021-10-31} \ganttnewline

			\ganttgroup{Prototype}{2021-10-04}{2021-10-31}
			\ganttgroup{Demonstration}{2021-11-03}{2021-12-05}
			\ganttnewline

			\ganttbar{Simulation}{2021-10-04}{2021-10-18}
			\ganttbar{Hardware}{2021-10-18}{2021-10-31}
			\ganttmilestone{Working SDR TX -- RX}{2021-11-01}

			\ganttbar{Fading model}{2021-11-03}{2021-11-19}
			\ganttbar{Measurements}{2021-11-20}{2021-12-05}
			\ganttmilestone{TX -- RX with Fading}{2021-12-06}
			\ganttnewline

			\ganttgroup{Documentation}{2021-09-21}{2021-12-20} \ganttnewline
			\ganttbar{Project Plan}{2021-09-21}{2021-10-04}
			\ganttmilestone{Submission}{2021-10-05}
			\ganttbar{Slides}{2021-12-06}{2021-12-22}
			\ganttmilestone{Presentation}{2021-12-23}
			\ganttnewline

			\ganttbar{Introduction and Task Desc.}{2021-09-27}{2021-10-24}
			\ganttbar{Methods (SW and HW)}{2021-10-27}{2021-11-16}
			\ganttbar{Results}{2021-11-20}{2021-12-11}
			\ganttmilestone{Ready}{2021-12-12}
			\ganttbar{Review}{2021-12-14}{2021-12-23}
			\ganttmilestone{Submission}{2021-12-24}
			\ganttnewline

			\ganttbar{Theory}{2021-10-07}{2021-12-05}
		\end{ganttchart}
	\end{figure}
\end{landscape}
\restoregeometry % Wieder die alten Ränder
\end{document}
